\section{Kết quả trên TAWOS}

\begin{frame}{Kết quả chính: giả thuyết tại mốc 50\%}
  \small
  \renewcommand{\arraystretch}{1.35}
  \begin{tabular}{l|r|r|r|l|l}
    \textbf{Thiết kế} & \textbf{Tĩnh} & \textbf{Động} & \textbf{Chênh lệch} & \textbf{CI 95\%} & \textbf{Đơn vị} \\
    \hline
    Temporal & 44,46\% & 53,36\% & \alert{+8,90} & [6,38; 11,70] & 336 sprint \\
    Cross-project & 43,56\% & 53,52\% & \alert{+9,96} & [7,53; 12,82] & 14 dự án \\
  \end{tabular}

  \vspace{0.3cm}
  \footnotesize
  \begin{itemize}
    \setlength{\itemsep}{2pt}
    \item \alert{Đọc:} với khoảng 20\% cảnh báo mỗi sprint, mô hình động bắt
          được trung bình 53\% số issue sẽ trễ, mô hình tĩnh bắt được 44\%.
          CI không chứa 0, nên \alert{giả thuyết chính được ủng hộ}.
    \item \alert{Cùng 850 cảnh báo (temporal):} phát hiện 682 so với 605 issue
          (trên tổng 1.883); precision gộp 80,24\% so với 71,18\%; cảnh báo sai
          168 so với 245 (0,42 so với 0,62 mỗi sprint).
  \end{itemize}
\end{frame}

\begin{frame}{Chính xác hơn nhưng muộn hơn}
  \small
  \renewcommand{\arraystretch}{1.35}
  \begin{tabular}{L{0.28\textwidth}|R{0.15\textwidth}|R{0.13\textwidth}|R{0.13\textwidth}|R{0.13\textwidth}}
    \textbf{Mốc} & \textbf{0\%} & \textbf{25\%} & \textbf{50\%} & \textbf{75\%} \\
    \hline
    Recall động & 45,12\% & 50,22\% & 53,36\% & 54,55\% \\
    Chênh lệch (điểm \%) & +0,66\textsuperscript{*} & +5,76 & +8,90 & +10,09 \\
    Lead time TB & $\sim$12,6 ngày & $\sim$9,5 ngày & $\sim$6,3 ngày & $\sim$3,1 ngày \\
  \end{tabular}

  \vspace{0.1cm}
  \scriptsize \textsuperscript{*} CI 95\% [$-$0,89; 2,04] chứa 0. Recall tĩnh = 44,46\% ở mọi mốc.

  \vspace{0.2cm}
  \footnotesize
  \begin{itemize}
    \setlength{\itemsep}{2pt}
    \item Ở mốc 0\% chưa có lịch sử nên hai mô hình ngang nhau. Đây là kết quả
          kỳ vọng và cho thấy không có rò rỉ thông tin.
    \item Mốc muộn hơn cho chênh lệch lớn hơn nhưng ít thời gian xử lý hơn. Đề tài
          \alert{không chọn một mốc ``tối ưu''} duy nhất.
  \end{itemize}
\end{frame}

\begin{frame}{Kiểm tra độ vững: phải đọc kỹ}
  \scriptsize
  \setlength{\tabcolsep}{3pt}
  \renewcommand{\arraystretch}{1.2}
  \begin{tabular}{L{0.42\textwidth}|L{0.2\textwidth}|L{0.2\textwidth}}
    \textbf{Chênh lệch động -- tĩnh tại 50\% (điểm \%)} & \textbf{Temporal} & \textbf{Cross-project} \\
    \hline
    Chính: $\lceil 0{,}2n\rceil$ & +8,90 [6,38; 11,70] & +9,96 [7,53; 12,82] \\
    Làm tròn xuống $\lfloor 0{,}2n\rfloor$ & +3,82 [2,20; 5,67] & +6,09 [4,34; 8,09] \\
    \alert{Chỉ issue còn mở} & \alert{+2,59 [0,68; 4,69]} & \alert{+2,03 [0,73; 3,58]} \\
    Chỉ tính Closed là hoàn thành & +1,87 [0,83; 3,04] & +1,95 [0,32; 3,81] \\
  \end{tabular}

  \vspace{0.15cm}
  \begin{alertblock}{Ba lưu ý}
    \begin{enumerate}
      \setlength{\itemsep}{1pt}
      \item Phần lớn lợi thế đến từ việc \alert{nhận ra issue đã Done}:
            \texttt{dynamic\_is\_done} chiếm 27,41\% / 56,44\% importance.
      \item Ngân sách thực tế \alert{không phải 20\%}: 207/397 sprint có dưới 5 issue,
            nên tỷ lệ cảnh báo thực tế là 41,54\% (macro) / 25,81\% (micro).
      \item Mọi chênh lệch vẫn dương và có ý nghĩa, nhưng nhỏ hơn nhiều so với con số chính.
    \end{enumerate}
  \end{alertblock}
\end{frame}

\begin{frame}{Chất lượng xác suất (mốc 50\%)}
  \footnotesize
  \setlength{\tabcolsep}{3pt}
  \renewcommand{\arraystretch}{1.35}
  \begin{tabular}{L{0.15\textwidth}|L{0.07\textwidth}|R{0.1\textwidth}|R{0.1\textwidth}|R{0.11\textwidth}|L{0.17\textwidth}|L{0.15\textwidth}}
    & & \textbf{AP TB dự án} & \textbf{Brier thô} & \textbf{Brier HC}
      & \textbf{Intercept} & \textbf{Slope} \\
    \hline
    Temporal & tĩnh & 0,6571 & 0,2188 & 0,1894 & 0,49 $\to$ 0,03 & 0,66 $\to$ 0,88 \\
             & động & 0,8128 & 0,1279 & 0,1137 & 0,57 $\to$ 0,09 & 0,87 $\to$ 1,03 \\
    \hline
    Cross-project & tĩnh & 0,6359 & 0,2257 & 0,2241 & 0,17 $\to$ $-$0,09 & 0,67 $\to$ 0,74 \\
                  & động & 0,8035 & 0,1374 & \textcolor{templateRed}{0,1381} & 0,20 $\to$ $-$0,004 & 0,89 $\to$ 0,92 \\
  \end{tabular}

  \vspace{0.1cm}
  \scriptsize HC: sau hiệu chuẩn. Lý tưởng: intercept 0, slope 1. Đoán ngẫu
  nhiên: AP $\approx$ 0,47--0,57; Brier $\approx$ 0,25.

  \vspace{0.15cm}
  \footnotesize
  \begin{itemize}
    \setlength{\itemsep}{2pt}
    \item Mô hình động \alert{xếp hạng tốt hơn rõ rệt} (AP $\sim$0,65 $\to$ $\sim$0,81).
    \item Cross-project: Brier \alert{không cải thiện}, nên cần hiệu chuẩn lại
          trên dữ liệu Agentick.
  \end{itemize}
\end{frame}

\begin{frame}{So với quy tắc đơn giản và giữa các dự án (khám phá)}
  \small
  \renewcommand{\arraystretch}{1.35}
  \begin{tabular}{l|l|l}
    \textbf{Chênh lệch tại 50\% (điểm \%)} & \textbf{Temporal} & \textbf{Cross-project} \\
    \hline
    Động -- quy tắc nghiệp vụ & \alert{+2,13 [$-$0,33; 4,60]} & +3,83 [2,10; 5,73] \\
    Động -- không cohort & $-$0,87 [$-$2,26; 0,29] & +0,72 [0,05; 1,42] \\
    Tuần tự: động -- tĩnh & +1,30 [$-$0,46; 3,21] & +1,88 [0,60; 3,14] \\
  \end{tabular}

  \vspace{0.25cm}
  \footnotesize
  \begin{itemize}
    \setlength{\itemsep}{2pt}
    \item Trên sprint tương lai, \alert{chưa chứng minh ML hơn quy tắc}
          (recall 51,23\% so với 53,36\%; CI chứa 0).
    \item Cảnh báo tuần tự: recall 51,56\%, precision 77,67\%, lead time
          $\sim$8,3 ngày, so với $\sim$6,3 ngày khi chỉ cảnh báo một lần ở 50\%.
    \item Giữa các dự án: chênh lệch động -- tĩnh \alert{dương ở cả 14 dự án} nhưng
          dao động mạnh (temporal +2,6 đến +35 điểm), nên không suy ra được mức
          cải thiện cho một nhóm cụ thể.
  \end{itemize}
\end{frame}

\begin{frame}{Phân bổ cảnh báo theo thời gian (khám phá)}
  \footnotesize
  \textbf{Câu hỏi:} giữ tổng $K$ cảnh báo, chỉ xét issue còn mở. Phát sớm một
  phần (\emph{quota}) có tốt hơn phát hết ở mốc 50\% (\emph{midpoint}) không?

  \vspace{0.15cm}
  \renewcommand{\arraystretch}{1.3}
  \begin{tabular}{L{0.32\textwidth}|R{0.16\textwidth}|R{0.16\textwidth}}
    $K=2$, temporal & \textbf{Quota} & \textbf{Midpoint} \\
    \hline
    Recall phát hiện sớm & 58,77\% & 59,20\% \\
    Precision & 77,31\% & 79,70\% \\
    Lead time & 8,18 ngày & 6,29 ngày \\
  \end{tabular}

  \vspace{0.2cm}
  \begin{itemize}
    \setlength{\itemsep}{2pt}
    \item Chênh lệch recall (quota -- midpoint): $-$0,43 [$-$1,36; 0,55] (temporal); $-$1,08 [$-$1,88; $-$0,35] (cross-project).
    \item \alert{Kết luận:} sớm hơn $\sim$1,9 ngày nhưng \alert{không tăng recall}.
          $K=2$ là giá trị được khảo sát, \alert{không phải} giá trị tối ưu đã chứng minh.
  \end{itemize}
\end{frame}
